% Author: Shuyu Fan, MSc
% Date: January 21, 2026

\subsection{Synthetic Patient Generation}

The test harness generates synthetic patient profiles for evaluating clinical decision support systems. The main function is \texttt{gen\_demo\_2()}, which outputs patient records with demographics and physical measurements.

\subsubsection{Input Parameters}

The function accepts:
\begin{itemize}
    \item \textbf{country}: Country for localized names and regional context
    \item \textbf{sample\_size}: Number of patients to generate (default: 1000)
    \item \textbf{area}: Region or state/province code
    \item \textbf{postal\_code}: Postal/ZIP code prefix
    \item \textbf{range\_day}: Date range for birth dates, which controls age distribution
    \item \textbf{reference\_date}: Date for age calculation (default: current date)
    \item \textbf{age\_cutoff}: Minimum age for female pregnancy status (default: 11 years)
    \item \textbf{p\_preg}: Pregnancy probability for eligible females (default: 0.1)
\end{itemize}

\subsubsection{Demographics Generation}

Each patient record includes:
\begin{itemize}
    \item \textbf{Unique Identifier (UID)}: Composite key in the format \texttt{Country\_Area\_PostalCode\_RID}, where RID is a sequential identifier starting from 10000
    \item \textbf{Gender}: Randomly assigned with 50/50 probability (1 = Male, 2 = Female)
    \item \textbf{Date of Birth}: Uniformly sampled from the specified date range
    \item \textbf{Age}: Stored in multiple formats: years, months, and combined (e.g., ``2Y5M'' for 2 years 5 months)
\end{itemize}
For a complete list of output fields, see Table~\ref{tab:patient_schema}.

\subsubsection{Name Localization}

Names come from the Faker library using locale-specific generators. The system has 52 country-locale mappings covering:
\begin{itemize}
    \item European languages: German, French, Spanish, Italian, Portuguese, Dutch, Polish, and others
    \item Asian languages: Chinese, Japanese, Korean, Thai, Vietnamese, Indonesian
    \item Middle Eastern languages: Hebrew, Farsi, and Arabic regional variants
    \item South Asian languages: Hindi, Tamil, Nepali, Bengali
    \item Others: Russian, Ukrainian, Turkish, Greek
\end{itemize}
The generator picks male or female names based on the assigned gender.

\subsubsection{Age Categories}

Patients are classified into age groups:
\begin{itemize}
    \item $<$ 1 year (infant)
    \item $<$ 2 years (toddler)
    \item $<$ 5 years (preschool)
    \item $<$ 10 years (school age)
    \item $<$ 14 years (early adolescent)
    \item $<$ 18 years (adolescent)
    \item $\geq$ 18 years (adult)
\end{itemize}
A separate field marks each patient as ``Adult'' or ``Child'' using the 18-year cutoff.

\subsubsection{Physical Attributes}

Height, weight, and BMI come from WHO/CDC pediatric growth reference data:
\begin{itemize}
    \item \textbf{Data sources}: WHO Child Growth Standards\footnote{\url{https://www.who.int/tools/child-growth-standards/standards}} (0--10 years) and CDC Growth Charts\footnote{\url{https://www.cdc.gov/growthcharts/percentile_data_files.htm}} (2--20 years):
    \begin{itemize}
        \item Infant length (0--24 months): \texttt{lenageinf.csv}
        \item Child/adolescent height (2--20 years): \texttt{statage.csv}
        \item Infant weight (0--24 months): \texttt{wtageinf.csv}
        \item Child/adolescent weight (2--20 years): \texttt{wtage.csv}
    \end{itemize}
    \item \textbf{Sampling}: Values are drawn from a uniform distribution between the 3rd and 97th percentiles for the patient's age and sex
    \item \textbf{BMI}: Calculated as $\mathrm{BMI} = \frac{\mathrm{weight~(kg)}}{\mathrm{height~(m)}^2}$
\end{itemize}

\subsubsection{Pregnancy Status}

For females above the age cutoff, pregnancy status is assigned randomly with probability \texttt{p\_preg}. Males and younger females get ``NA''.
\begin{table}[h]
\centering
\begin{tabular}{ll}
\hline
\textbf{Field} & \textbf{Description} \\
\hline
COUNTRY & Country of residence \\
AREA & Region/state/province \\
POSTAL\_CODE & Postal code prefix \\
RID & Record identifier \\
UID & Unique patient identifier \\
NAME & Localized patient name \\
GENDER & 1 (Male) or 2 (Female) \\
DOB & Date of birth \\
AGE & Age in years \\
AGE\_IN\_STR & Formatted age string (e.g., ``2Y5M'') \\
AGE\_CATEGORY & Age group \\
ADULT\_CHILD & Adult or Child \\
HeightCM & Height in centimeters \\
WeightKG & Weight in kilograms \\
BMI & Body mass index \\
PREGNANT & Yes/No/NA \\
\hline
\end{tabular}
\caption{Synthetic patient output schema}
\label{tab:patient_schema}
\end{table}

\subsubsection{Clinical Context Assignment}

After generating demographics, the system assigns conditions and symptoms from the knowledge graph:
\begin{itemize}
    \item \textbf{Conditions}: Picked from the disease dictionary for the evaluation package (e.g., HIV, immunization)
    \item \textbf{Symptoms}: 1--3 symptoms selected from the condition's list, weighted by clinical importance
    \item \textbf{Age-specific presentation}: Selection accounts for pediatric vs. adult symptom patterns
\end{itemize}
